% Authored. Tables and numbers come from generated/.

\section{Data and Harmonisation}
\label{sec:data}

\subsection{Sources}

The canonical balance panel takes \PanelStart--1994 from the Banco de Portugal /
INE historical long series \citep{bportugal_series} and 1995--\PanelEnd\ from INE
national accounts distributed through PORDATA \citep{pordata_2785}. Detailed
revenue, expenditure, interest, investment, Maastricht debt and stock-flow data for
the modern period come from the ESA 2010 annual workbooks published by the
Portuguese Public Finance Council \citep{cfp_data}, which also serve as an
independent cross-check on the balance series. Table~\ref{tab:sources} records each
file, the coverage taken from it, its vintage and a digest of the retained copy.

Every workbook is parsed programmatically from the bundled file. No value in this
paper was transcribed by hand, including the values in this sentence: they are
macros written by the pipeline from the persisted artefacts.

\begin{table}[htbp]
\centering
\caption{Bundled sources, the coverage taken from each, the vintage of the retained file and the first twelve hexadecimal digits of its SHA-256 digest.}
\label{tab:sources}
\small
\fitwidth{%
\begin{tabular}{@{}lllll@{}}
\toprule
Source & Institution & Coverage used & Vintage & SHA-256 prefix \\
\midrule
Banco Portugal Long Series & Banco de Portugal / INE & 1977-1995 for the historical segment used here & 2023-12 & d3570ea90f78 \\
PORDATA Balance By Level & INE via PORDATA & 1995-2025 & 2026-04 & 2a5d54649041 \\
CFP Annual General Government & Portuguese Public Finance Council (CFP) & 1995-2025 & 2026-04-15 & 8f6b5a2916a4 \\
CFP Annual Subsectors & Portuguese Public Finance Council (CFP) & 2000-2025 & 2026-04-15 & 4d9e1b07cd38 \\
CFP Social Security Detail & Portuguese Public Finance Council (CFP) & 2019-2025 systems; detailed 2024-2025 tables & 2026-04 & d9f28e424ab6 \\
\bottomrule
\end{tabular}%
}
\end{table}


\subsection{Vintage}

A fiscal series has a vintage as well as a coverage, and the two are different
things. The modern national-accounts and CFP files used here are the April 2026
releases. Within them, \ProvisionalYears\ are flagged \textbf{provisional} by the
publisher, and that flag is carried through the harmonisation into the canonical
panel rather than dropped.

The distinction is not pedantry in this application. The provisional years are the
years a reader is most likely to be interested in, and they are the years most
likely to be restated. Any statement below about \ProvisionalYears\ should be read
as a statement about the April 2026 vintage of those years.

\subsection{Two statistical regimes, not one series}

The splice at \textbf{1995} separates two source families built on different
methodologies. We retain it explicitly. Nothing is chained, smoothed, interpolated
or calibrated across it, and two consequences are enforced throughout the paper:

\begin{enumerate}
\item no statistic whose value depends on the \emph{level} of the balance --- a
mean, a segment mean, a regression intercept --- is computed across the boundary;
\item no annual change is computed through it, so the 1994-to-1995 change, which
would mix a vintage revision with an economic movement, appears nowhere in the
results.
\end{enumerate}

1995 exists in both source families, so rather than choosing silently we keep both
vintages and report the revision between them. The revisions are small in level
terms but non-zero for every subsector, which is the direct evidence that the two
regimes are not interchangeable at the level of a single year.

\subsubsection*{Signs are more robust to the splice than magnitudes --- checked, not assumed}

A natural claim at this point is that a balance's sign is insensitive to the vintage
while its level is not. Stated generally that is false: a methodological revision can
move a balance close to zero across it, turning a small surplus into a small deficit.

What can be checked is whether it happens here, and in the overlaps this panel
retains it does not. Across the \OverlapSignTotal\ balances observed in both 1995
vintages, the sign is the same in \OverlapSignAgreements. Across the
\SourceSignComparisons\ subsector-years in which the PORDATA bridge and the CFP
workbook both report a figure, the sign agrees in \SourceSignAgreements. We therefore
treat sign classifications as empirically invariant across the source overlaps
available here, and as less exposed to the splice than magnitudes --- which is a claim
about this panel rather than a general property of fiscal statistics.

\subsection{An explicit gap}

Detailed revenue and expenditure components are available for General Government
continuously from \PanelStart. For the three subsectors they exist for
\PanelStart--1995 and 2000--\PanelEnd\ only. The \GapYears\ intervening years are
absent from both source families and are therefore absent here: they are not
imputed, no statistic is computed across them, and every figure drawn from the
account panel breaks its line rather than joining 1995 to 2000.

The gap affects the detailed components only. The canonical B.9 panel has an
observation for every one of the \PanelYears\ years and for every subsector, which
is why the composition results in Section~\ref{sec:results} cover the full window
while the revenue--expenditure results do not. The detailed account panel holds
\AccountObservations\ sector-year observations against the
$4 \times \PanelYears$ a complete panel would have, and
Table~\ref{tab:coverage} states the two coverages separately: reporting a single
observation count for a panel that has two would be misleading.

\begin{table}[htbp]
\centering
\caption{Coverage of the detailed account panel. The canonical balance panel is complete; only the detailed components carry the 1996--1999 gap.}
\label{tab:coverage}
\small
\fitwidth{%
\begin{tabular}{@{}lrrr@{}}
\toprule
Subsector & First year & Last year & Observations \\
\midrule
General Government & 1977 & 2025 & 49 \\
Central Government & 1977 & 2025 & 45 \\
Regional and Local & 1977 & 2025 & 45 \\
Social Security Funds & 1977 & 2025 & 45 \\
\bottomrule
\end{tabular}%
}
\end{table}


\subsection{Identity closure is not source agreement}

Two different validation questions are easy to conflate, and conflating them
overstates what the checks establish.

An \textbf{identity} check asks whether an extraction is arithmetically
self-consistent: whether B.9 equals the sum of its subsectors, whether revenue
minus expenditure reproduces the recorded balance, whether the change in debt
reconciles with the balance and the stock-flow adjustment. Such a check can close
to numerical precision while both sources are wrong in the same way, because it
never consults a second source.

A \textbf{source-agreement} check asks whether two independently published sources
report the same number for the same year. Identity closure cannot establish it.

Table~\ref{tab:validation} reports both in one unit. The identities close to
\MaxClosureResidual\ million euro, the rounding of the published series. The largest
disagreement between the PORDATA bridge and the CFP workbook is
\MaxSourceDisagreement\ million euro, on Central Government in
\SourceDisagreementWhere. That gap is nearly two orders of magnitude larger than any
identity residual, which is the concrete reason for reporting the two separately:
a paper that reported only the identity residuals would imply a precision the
underlying statistics do not have.

\begin{table}[htbp]
\centering
\caption{Identity closure and source agreement, in one unit. The two answer different questions and neither substitutes for the other.}
\label{tab:validation}
\small
\fitwidth{%
\begin{tabular}{@{}llrrl@{}}
\toprule
Check & Quantity compared & N & Largest absolute difference (M EUR) & Where \\
\midrule
Accounting identity & Aggregate balance minus the sum of its three subsectors & 49 & 1.000 & 2004 \\
Accounting identity & Revenue minus expenditure minus the recorded balance & 184 & 0.000 & Central Government, 1994 \\
Accounting identity & Debt change against minus the balance plus the stock-flow adjustment & 108 & 0.000 & 2000 \\
Source agreement & PORDATA bridge against the CFP workbook, General Government & 31 & 0.498 & 1995 \\
Source agreement & PORDATA bridge against the CFP workbook, Central Government & 26 & 66.811 & 2002 \\
Source agreement & PORDATA bridge against the CFP workbook, Regional and Local & 26 & 0.492 & 2002 \\
Source agreement & PORDATA bridge against the CFP workbook, Social Security Funds & 26 & 0.470 & 2019 \\
Vintage revision & 1995 modern vintage against the historical vintage, four balances & 4 & 72.465 & 1995 \\
\bottomrule
\end{tabular}%
}
\end{table}

